% Draft-specific packages and working-draft markers for draft-initial.
\usepackage{placeins}
\usepackage{pgfplots}
\pgfplotsset{compat=1.17}
\usetikzlibrary{arrows.meta,positioning,fit,calc,shapes.geometric,decorations.pathreplacing,patterns}
\input{"../draft-initial/glossary.tex"}

\hypersetup{pdftitle={ELL305 Term Paper - draft-initial (working draft)}}

% Article class has no chapters: Parts are \part, chapters are \section.
\titleformat{\part}[display]
  {\centering\color{draftblue}\Huge\bfseries}{Part \thepart}{0.6em}{}
\setcounter{secnumdepth}{3}
\setcounter{tocdepth}{2}

% Working-draft markers. Every marker must be resolved or removed before
% this material is used in any submitted milestone.
\definecolor{todored}{HTML}{B0222B}
\definecolor{verifyamber}{HTML}{8A5A00}
\newcommand{\todo}[1]{{\color{todored}\textbf{[TODO:} #1\textbf{]}}}
\newcommand{\verify}[1]{{\color{verifyamber}\textbf{[VERIFY:} #1\textbf{]}}}
\newcommand{\partbudget}[2]{%
  \begin{center}\small\itshape\color{draftgray}
  Target length: about #1 pages. Draft status: #2.
  \end{center}}
\newcommand{\plannedcontent}[1]{%
  \par\noindent\fcolorbox{draftgray}{white}{%
    \begin{minipage}{0.96\textwidth}\small
    \textbf{Planned content (not yet written).} #1
    \end{minipage}}\par}
